%! Unit 1: Algebraic Expressions and Factoring — Summative Assessment — Practice Test --- Study Copy (blank)
%
% Blueprint (100 pts): Part A vocabulary 8x1, Part B multiple choice 6x2,
% Part C short answer 10x5 (ONE ITEM PER LESSON, in lesson order), Part D
% extended response 2x15. The actual test is the parallel form: same blueprint,
% different numbers, reshuffled vocabulary letters.
%
% WORK RULE: every computation is followed by a \begin{work} block authored
% byte-identically here and in ../../test_keys/practice_test_key/main.tex. Under
% -boxes the block ships as a \vphantom, so the student gets exactly the room the
% answer needs; under -key the same block prints in keyred. That is what keeps
% this test and its key the same number of pages — which matters because
% shared/unit.mk swaps sample_test_key in for sample_test at the tail of the key
% packet. \workrowsep is the handwriting room, and it moves both files together.
\documentclass[10pt]{article}
\usepackage{atda-article}
\usepackage{atda-boxes}
\newcommand{\TallMath}[1]{$\displaystyle #1\rule[-1.4em]{0pt}{3.2em}$}
\setlength{\workrowsep}{8pt}

% Part-divider strip.
% BOXGUARD: the guard lives inside the macro, so a part header can never be
% stranded alone at the foot of a page, in the blank or the key — both files use
% this identical definition.
\newcommand{\parthead}[1]{%
  \par\boxguard[9]%
  \vspace{6pt}\noindent
  \begin{headlinebox}{royal}{\color{white}\bfseries #1}\end{headlinebox}%
  \vspace{4pt}}

\begin{document}

\pageheader{Unit 1: Algebraic Expressions and Factoring}{Practice Test --- Study Copy}
\namedateperiod

\vspace{0.10in}
\boxguard[8]
\begin{remindbox}
\small
\textbf{This is a practice test.} It mirrors the real Unit 1 test in length, format,
and the ideas it covers, but the numbers and contexts differ. Work every problem
with no notes, then check your answers against the key.
\end{remindbox}

% ======================================================================
\parthead{Part A --- Vocabulary \quad (8 items, 1 point each --- 8 points)}

\noindent\small
Write the letter of the best definition on the line beside each term. \textbf{Two definitions
will not be used.}

\vspace{5pt}
\noindent
\begin{tabularx}{\linewidth}{@{} p{0.95cm} X p{0.95cm} X @{}}
\blank{0.7cm} & \textbf{1.} Greatest common factor (GCF)
  & \blank{0.7cm} & \textbf{5.} Zero product property \\[4pt]
\blank{0.7cm} & \textbf{2.} Prime polynomial
  & \blank{0.7cm} & \textbf{6.} Double root \\[4pt]
\blank{0.7cm} & \textbf{3.} Perfect square trinomial
  & \blank{0.7cm} & \textbf{7.} Excluded value \\[4pt]
\blank{0.7cm} & \textbf{4.} Difference of squares
  & \blank{0.7cm} & \textbf{8.} Rational expression \\
\end{tabularx}

\vspace{6pt}
\begin{enumerate}[label=(\Alph*), leftmargin=2.4em, itemsep=1pt, topsep=2pt, font=\bfseries]
  \small
  \item A quotient of two polynomials.
  \item A binomial of the form $a^2-b^2$; it factors as $(a+b)(a-b)$.
  \item The largest factor shared by every term of a polynomial.
  \item A number that makes a denominator zero, so the expression is undefined there.
  \item An expression written as a product of two or more expressions.
  \item A polynomial with integer coefficients that cannot be factored any further over the integers.
  \item If a product equals zero, then at least one of its factors equals zero.
  \item A solution that appears twice because the factor producing it is squared.
  \item The exponent that tells how many times a base is used as a factor.
  \item A trinomial of the form $a^2 \pm 2ab + b^2$; it factors as $(a \pm b)^2$.
\end{enumerate}

% ======================================================================
\parthead{Part B --- Multiple Choice \quad (6 items, 2 points each --- 12 points)}

\noindent\small Circle the best answer.

\begin{enumerate}[leftmargin=1.8em, itemsep=6pt, topsep=5pt]
\small

\item Which expression is equal to $(2m+3)^2$?
\begin{enumerate}[label=(\Alph*), leftmargin=2.2em, itemsep=1pt, topsep=2pt]
  \item $4m^2+9$
  \item $4m^2+6m+9$
  \item $4m^2+12m+9$
  \item $2m^2+12m+9$
\end{enumerate}

\item What is the value of $\sqrt{9+16}$?
\begin{enumerate}[label=(\Alph*), leftmargin=2.2em, itemsep=1pt, topsep=2pt]
  \item $7$
  \item $5$
  \item $25$
  \item $12$
\end{enumerate}

\item Which is the \textbf{complete} factorization of $6x^2+9x$?
\begin{enumerate}[label=(\Alph*), leftmargin=2.2em, itemsep=1pt, topsep=2pt]
  \item $x(6x+9)$
  \item $3(2x^2+3x)$
  \item $3x(2x+3)$
  \item $3x(2x^2+3)$
\end{enumerate}

\item Which polynomial is \textbf{prime} over the integers?
\begin{enumerate}[label=(\Alph*), leftmargin=2.2em, itemsep=1pt, topsep=2pt]
  \item $x^2-49$
  \item $x^2+49$
  \item $x^2-14x+49$
  \item $x^2+14x+49$
\end{enumerate}

\item How many solutions does $x^3+9x=0$ have, and of what type?
\begin{enumerate}[label=(\Alph*), leftmargin=2.2em, itemsep=1pt, topsep=2pt]
  \item $3$ solutions: $3$ real
  \item $3$ solutions: $1$ real and $2$ imaginary
  \item $1$ solution: $1$ real
  \item $2$ solutions: $2$ real
\end{enumerate}

\item Which simplification is \textbf{correct}?
\begin{enumerate}[label=(\Alph*), leftmargin=2.2em, itemsep=1pt, topsep=2pt]
  \item $\dfrac{x+5}{x+7} = \dfrac{5}{7}$
  \item $\dfrac{x+5}{5} = x$
  \item $\dfrac{5(x+5)}{5(x+7)} = \dfrac{x+5}{x+7}$
  \item $\dfrac{x^2+5}{x^2} = 5$
\end{enumerate}

\end{enumerate}

% ======================================================================
\parthead{Part C --- Short Answer \& Computation \quad (10 items, 5 points each --- 50 points)}

\noindent\small Show all work. An answer with no supporting work earns no credit.

\begin{enumerate}[leftmargin=1.8em, itemsep=4pt, topsep=5pt]
\small

\item Simplify completely. Write your answer with positive exponents only.
\ \TallMath{\frac{\left(3x^2y^3\right)^2 \cdot 2xy}{6x^3y^4}}
\begin{work}
  \frac{\left(3x^2y^3\right)^2 \cdot 2xy}{6x^3y^4}
      &= \frac{9x^4y^6 \cdot 2xy}{6x^3y^4} \\
      &= \frac{18x^5y^7}{6x^3y^4} \\
      &= 3x^2y^3
\end{work}

\item (a) Simplify $\sqrt{75x^3}$, where $x>0$. \quad (b) Write $\sqrt[3]{x^5}$ using a rational
exponent. \quad (c) Evaluate $16^{3/4}$.
\begin{work}
  \text{(a)} \quad \sqrt{75x^3} &= \sqrt{25x^2}\cdot\sqrt{3x} = 5x\sqrt{3x} \\
  \text{(b)} \quad \sqrt[3]{x^5} &= x^{5/3} \\
  \text{(c)} \quad 16^{3/4} &= \left(\sqrt[4]{16}\right)^3 = 2^3 = 8
\end{work}

\item Factor completely by grouping: \ $3x^3+12x^2+2x+8$
\begin{work}
  3x^3+12x^2+2x+8 &= 3x^2(x+4) + 2(x+4) \\
      &= (x+4)\left(3x^2+2\right)
\end{work}

\item Factor completely: \ $4x^2+11x+6$
\begin{work}
  ac = 4\cdot 6 &= 24, \quad \text{and } 3+8 = 11 \\
  4x^2+11x+6 &= 4x^2+8x+3x+6 \\
      &= 4x(x+2) + 3(x+2) \\
      &= (x+2)(4x+3)
\end{work}

\item Factor completely. \quad (a) $49x^2-64$ \qquad (b) $x^3-27$
\begin{work}
  \text{(a)} \quad 49x^2-64 &= (7x)^2 - 8^2 = (7x+8)(7x-8) \\
  \text{(b)} \quad x^3-27 &= x^3 - 3^3 = (x-3)\left(x^2+3x+9\right)
\end{work}

\item Solve by factoring: \ $2x^2 = 5x+12$
\begin{work}
  2x^2-5x-12 &= 0 \\
  (x-4)(2x+3) &= 0 \\
  x-4 = 0 \quad \text{or} \quad 2x+3 &= 0 \\
  x = 4 \quad \text{or} \quad x &= -\tfrac{3}{2}
\end{work}

\item Simplify, and state \textbf{every} excluded value.
\ \TallMath{\frac{x^2-16}{x^2-x-12}}
\begin{work}
  \frac{x^2-16}{x^2-x-12} &= \frac{(x+4)(x-4)}{(x-4)(x+3)} \\
      &= \frac{x+4}{x+3} \qquad (x \ne 4, \; x \ne -3)
\end{work}

\item Divide and simplify, and state \textbf{every} excluded value.
\ \TallMath{\frac{x^2-4}{x^2+6x+9} \div \frac{x+2}{x+3}}
\begin{work}
  \frac{(x+2)(x-2)}{(x+3)^2} \cdot \frac{x+3}{x+2}
      &= \frac{x-2}{x+3} \qquad (x \ne -3, \; x \ne -2)
\end{work}

\item A rectangular banner has area $A(x) = 2x^2+11x+12$ square feet and width $x+4$ feet.
(a) Find an expression for its length. (b) Find the length, width, and area when $x=5$.
\begin{work}
  2x^2+11x+12 &= (x+4)(2x+3) \\
  \text{(a)} \quad \text{length} &= 2x+3 \text{ feet} \\
  \text{(b) at } x=5: \quad \text{width} &= 9 \text{ ft}, \quad \text{length} = 13 \text{ ft} \\
  \text{area} &= 9 \cdot 13 = 117 \text{ square feet}
\end{work}

\item Factor completely. \quad (a) $5x^3-45x$ \qquad (b) $2x^4-32$
\begin{work}
  \text{(a)} \quad 5x^3-45x &= 5x\left(x^2-9\right) = 5x(x+3)(x-3) \\
  \text{(b)} \quad 2x^4-32 &= 2\left(x^4-16\right) \\
      &= 2\left(x^2+4\right)\left(x^2-4\right) \\
      &= 2\left(x^2+4\right)(x+2)(x-2)
\end{work}

\end{enumerate}

% ======================================================================
\parthead{Part D --- Extended Response \quad (2 items, 15 points each --- 30 points)}

\begin{enumerate}[leftmargin=1.8em, itemsep=8pt, topsep=5pt]
\small

\item \textbf{Verify, then judge a claim.}

\begin{enumerate}[label=(\alph*), leftmargin=1.8em, itemsep=4pt, topsep=4pt]
  \item Verify the identity $(x-4)\left(x^2+4x+16\right) = x^3-64$ by multiplying.
\begin{work}
  (x-4)\left(x^2+4x+16\right)
      &= x^3+4x^2+16x - 4x^2-16x-64 \\
      &= x^3-64
\end{work}

  \item A student writes: ``$\dfrac{x^3-64}{x-4}$ simplifies to $x^2+4x+16$, \emph{for every value
        of $x$}.'' Test both expressions at $x=5$ and at $x=4$.
\begin{work}
  \text{at } x=5: \quad \frac{125-64}{1} &= 61, \qquad 25+20+16 = 61 \\
  \text{at } x=4: \quad \frac{64-64}{0} &= \frac{0}{0} \text{ --- undefined} \\
  \text{but} \quad 16+16+16 &= 48
\end{work}

  \item Is the student's algebra correct? Is the student's \emph{claim} correct? Explain in one or
        two sentences, using the word \textbf{excluded}.

\writelines{3}
\end{enumerate}

\item \textbf{A rocket is launched from a platform.} Its height in feet after $t$ seconds is
$h(t) = -16t^2+64t+80$.

\begin{enumerate}[label=(\alph*), leftmargin=1.8em, itemsep=4pt, topsep=4pt]
  \item Factor $h(t)$ completely and find both values of $t$ for which $h(t)=0$.
\begin{work}
  h(t) &= -16t^2+64t+80 \\
      &= -16\left(t^2-4t-5\right) \\
      &= -16(t-5)(t+1) \\
  t &= 5 \quad \text{or} \quad t = -1
\end{work}

  \item Which of those two values describes when the rocket lands? What does the other value mean
        here, and why is it rejected?

\writelines{2}

  \item Solve $h(t)=80$. Interpret \textbf{both} solutions --- this time neither one is rejected.
\begin{work}
  -16t^2+64t+80 &= 80 \\
  -16t^2+64t &= 0 \\
  -16t(t-4) &= 0 \\
  t = 0 \quad \text{or} \quad t &= 4
\end{work}

\writelines{2}

  \item How many solutions does the equation $h(t)=0$ have, and of what type? How does the degree
        of $h(t)$ tell you this before you factor anything?

\writelines{2}
\end{enumerate}

\end{enumerate}

\end{document}
